% Tailored from reusable AI/ML resume for micro1 ML Engineer contractor role.
\documentclass[9pt,letterpaper]{extarticle}
\usepackage[left=0.68in,right=0.68in,top=0.45in,bottom=0.45in]{geometry}
\usepackage[T1]{fontenc}
\input{glyphtounicode}
\pdfgentounicode=1
\usepackage[utf8]{inputenc}
\usepackage{lmodern}
\usepackage{microtype}
\usepackage{pifont}
\usepackage{xcolor}
\usepackage{enumitem}
\usepackage{titlesec}
\usepackage[hidelinks]{hyperref}
\usepackage{parskip,ragged2e,needspace}
\pagestyle{empty}
\setlength{\parindent}{0pt}
\setlength{\parskip}{1pt}
\definecolor{accent}{HTML}{155B8A}
\titleformat{\section}{\large\bfseries\color{accent}}{}{0pt}{}[\vspace{0.5pt}\color{accent}\titlerule]
\titlespacing*{\section}{0pt}{3pt}{2pt}
\setlist[itemize]{leftmargin=1.15em,itemsep=0.7pt,topsep=1pt,parsep=0pt,partopsep=0pt}
\newcommand{\role}[4]{\Needspace{4\baselineskip}\textbf{#1}\hfill\textbf{#2}\\[-1pt]\textit{#3}\hfill\textit{#4}\par}
\newcommand{\project}[3]{\Needspace{5\baselineskip}\textbf{#1}\hfill\textit{#2 \textbar{} #3}\par}
\newcommand{\cert}[3]{\Needspace{4\baselineskip}\textbf{#1}\hfill\textit{#2}\par{\small #3\par}}
\newcommand{\pub}[3]{\Needspace{3\baselineskip}\textbf{#1}\hfill\textit{#2}\par{\small #3\par}}
\hypersetup{
  pdftitle={Waseem Raza, PhD - micro1 ML Engineer},
  pdfauthor={Waseem Raza}
}

\begin{document}
\begin{center}
{\LARGE\bfseries Waseem Raza, PhD}\\[2pt]
{\large\bfseries Machine Learning Engineer | Model Training, Inference \& Python Systems}\\[1pt]
{\small Deep Learning | Numerical Computing | Reproducible ML | Evaluation \& Performance}\\[3pt]
{\footnotesize New Jersey, United States \quad
\ding{37}\enspace \href{tel:+17328039703}{+1 732-803-9703} \quad
\ding{41}\enspace \href{mailto:waseem.placements@gmail.com}{waseem.placements@gmail.com} \textbar{}
\href{mailto:waseemraza844@gmail.com}{waseemraza844@gmail.com}\\[1pt]
\href{https://www.linkedin.com/in/waseem-raza-phd-9a737858/}{LinkedIn} \textbar{}
\href{https://github.com/WaseemRaza844}{GitHub} \textbar{}
\href{https://scholar.google.com/citations?user=d58Drr0AAAAJ&hl=en}{Google Scholar} \textbar{}
\href{https://waseemraza844.github.io/portfolio/}{Portfolio}}
\end{center}
\vspace{-5pt}

\section{Professional Summary}
Machine Learning Engineer and researcher with a PhD in Electrical and Computer Engineering and 7+ years of combined industry and research experience developing, training, validating, and debugging machine-learning systems. Strong hands-on background in Python, PyTorch, TensorFlow/Keras, scikit-learn, XGBoost, NumPy, SciPy, Pandas, and Hugging Face Transformers, with experience spanning neural networks, autoencoders, GANs, reinforcement learning, transfer learning, domain adaptation, classification, anomaly detection, time-series analysis, and multimodal data. Built reproducible experimentation and data pipelines across AT\&T, Futurewei, university research, and cloud analytics environments, including 100+ controlled ML experiments, 500K+ operational-log analysis, and production-oriented model validation under distribution shift.

\section{Technical Skills}
\textbf{Python \& Numerical Computing:} Python, NumPy, SciPy, Pandas, SQL, C/C++, Jupyter, vectorized data processing, statistical analysis, numerical experimentation, Linux, Bash/CLI, Git.\\
\textbf{ML \& Deep Learning:} PyTorch, TensorFlow, Keras, scikit-learn, XGBoost, Hugging Face Transformers; DNNs, CNNs, LSTMs, autoencoders, GANs, reinforcement learning, transfer learning, domain adaptation, classification, regression, anomaly detection, forecasting.\\
\textbf{Training, Inference \& Evaluation:} data preparation, feature engineering, train/validation/test workflows, cross-validation, model comparison, hyperparameter experimentation, batch inference, metric tracking, error analysis, failure-mode assessment, robustness testing, distribution-shift evaluation.\\
\textbf{ML Engineering \& Reproducibility:} reusable Python pipelines, experiment automation, deterministic/reproducible workflows, data-quality checks, model-validation harnesses, Git-based development, CI/CD, GitHub Actions, Docker, Kubernetes, OpenShift.\\
\textbf{Data \& Compute Platforms:} Azure Databricks, Apache Spark/PySpark, Synapse Analytics, Azure Functions, cloud data pipelines, ETL/ELT, structured/unstructured data processing, Google Cloud ML workflows.\\
\textbf{AI-Assisted Development:} coding-agent-assisted Python development, debugging, test generation, refactoring, repository iteration, and structured review of implementation correctness.

\section{Professional Experience}
\role{Machine Learning Engineer}{02/2025 -- Present}{AT\&T, contractor via Innovatix Inc.}{Bedminster, NJ}
\begin{itemize}
\item Own end-to-end development, testing, and operational integration of machine-learning diagnostics for critical E911 service-assurance workflows, improving automated root-cause classification accuracy from 20\% to 80\%.
\item Design domain-adaptive neural models and reusable preprocessing pipelines that improve generalization across heterogeneous operational domains and distribution shifts.
\item Build cross-domain model-validation, error-analysis, and failure-mode assessment workflows to identify data/model defects and verify readiness for operational deployment.
\item Translate network and public-safety requirements into measurable ML objectives, validated datasets, repeatable experiments, and production-oriented decision-support workflows.
\end{itemize}

\role{Data Scientist and AI/ML Engineer}{09/2024 -- 12/2024}{Impac Exploration Services}{Weatherford, OK}
\begin{itemize}
\item Built machine-learning and LLM-assisted analytical workflows for large geochemical datasets using Python and Azure cloud services.
\item Developed scalable ingestion, cleansing, feature-preparation, transformation, and validation pipelines with Azure Databricks, Azure Functions, and Synapse Analytics.
\item Automated ETL workflows, reducing data-preparation time by 40\%; implemented cross-validation, data-quality, and consistency checks that improved reliability of analytical outputs by 20\%.
\end{itemize}

\newpage
\role{Graduate Research Assistant}{09/2019 -- 07/2025}{AI4Networks, University of Oklahoma}{Tulsa, OK}
\begin{itemize}
\item Designed reusable Python simulation and data-generation platforms for large-scale ML experimentation, including a 3GPP-compliant network simulator and AI-enabled 5G testbed.
\item Built and evaluated transfer-learning, adversarial-learning, GAN, autoencoder, and reinforcement-learning pipelines for sparse-data and cross-domain modeling.
\item Developed robustness-evaluation workflows for distribution shifts, missing data, and limited labeled samples; led 100+ experiments from problem formulation through quantitative validation.
\item Developed scalable self-healing and outage-management frameworks that improved modeled network resilience by 30\%.
\end{itemize}

\role{Advanced Wireless Technology Development Intern}{09/2022 -- 01/2023}{Futurewei Technologies}{Schaumburg, IL}
\begin{itemize}
\item Improved autoencoder-based channel-state-information compression by 25\% using transfer learning and systematic experimentation on large wireless datasets.
\item Conducted 50+ experiments to validate model accuracy, robustness, and generalization under changing data conditions; translated results into recommendations for next-generation communication systems.
\end{itemize}

\role{Advanced-Degrees Intern}{06/2024 -- 08/2024}{AT\&T Labs}{Austin, TX}
\begin{itemize}
\item Analyzed 500K+ LTE/5G-NR diagnostic logs and applied ML methods to detect anomalies, reducing diagnostic time by 20\%.
\item Improved data-driven decision algorithms through structured experimentation, model evaluation, and validation over large operational datasets.
\end{itemize}

\section{Education}
\textbf{PhD, Electrical and Computer Engineering}\hfill\textit{09/2019 -- 07/2025}\\
University of Oklahoma, Norman, Oklahoma, USA\\
\textbf{MSc, Telecommunication Engineering}\hfill\textit{08/2014 -- 08/2016}\\
University of Engineering and Technology, Taxila, Pakistan\\
\textbf{BSc, Telecommunication Engineering}\hfill\textit{11/2010 -- 06/2014}\\
University of Engineering and Technology, Taxila, Pakistan

\section{Selected ML Engineering Projects}
\project{Wearable Intelligence for Health Monitoring}{Research Project}{2024--2026}
Developed deep-learning workflows for continuous recognition of medication-taking and dietary activities from smartwatch accelerometer and gyroscope streams. Built preprocessing/windowing pipelines, trained CNN/LSTM-style models, compared sequence lengths and sampling configurations, and evaluated cross-user generalization with recall-oriented metrics.\\[-1pt]
{\small\textit{Tools \& methods:} Python, deep learning, CNN/LSTM, time-series segmentation, feature engineering, smartwatch IMU, train/test evaluation, cross-user validation, confusion matrices, model comparison.}

\project{Propagation Models \& Digital Twins}{Research Project}{2024}
Developed robust data-driven propagation models for wireless digital twins using domain-informed neural pretraining, GAN-based data augmentation, and scarce measured data. Evaluated synthetic-data fidelity and downstream regression performance, including robustness under domain shift and limited training data.\\[-1pt]
{\small\textit{Tools \& methods:} Python, deep neural networks, GANs, transfer/domain adaptation, synthetic data, regression, RMSE/R\textsuperscript{2}, statistical validation, 3GPP-informed modeling.}

\project{SyntheticNET: 3GPP-Compliant Network Simulation}{Research Project}{2020--2024}
Developed a Python-based simulation and synthetic-data environment for repeatable large-scale experiments when real measurements are sparse or expensive. Generated controlled mobility, propagation, and network-state datasets used to train, test, and compare ML approaches across consistent scenarios.\\[-1pt]
{\small\textit{Tools \& methods:} Python, numerical simulation, synthetic data generation, 3GPP models, experiment automation, data pipelines, reproducibility, mobility/propagation modeling.}

\project{Autoencoder CSI Compression \& Transfer Learning}{Industry Research Project}{2022--2023}
Developed and evaluated autoencoder-based channel-state-information compression workflows for large wireless datasets, using transfer learning to improve reconstruction/compression performance by 25\%. Ran 50+ controlled experiments across changing data conditions to quantify accuracy, robustness, and generalization before translating results into engineering recommendations.\\[-1pt]
{\small\textit{Tools \& methods:} Python, autoencoders, deep neural networks, transfer learning, representation learning, model compression for signal representations, train/validation experiments, robustness testing, performance benchmarking.}

\project{Financial ML and Forecasting}{Portfolio Project}{2024--2025}
Developed structured-data ML workflows for credit-risk and loan-default classification alongside time-series forecasting experiments. Implemented data preparation, feature engineering, model comparison, cross-validation, and predictive-performance analysis over finance-oriented datasets.\\[-1pt]
{\small\textit{Tools \& methods:} Python, scikit-learn, classification, time-series forecasting, feature engineering, cross-validation, model evaluation, predictive analytics.}

\project{Cloud Computer Vision}{Guided Project}{2025}
Implemented image-classification, object-detection, and vehicle-counting workflows with managed cloud models and Python inference pipelines. Validated preprocessing, prediction outputs, class/count aggregation, and model behavior across visual inputs.\\[-1pt]
{\small\textit{Tools \& methods:} Python, TensorFlow concepts, computer vision, object detection, image classification, Google Cloud AutoML, preprocessing, inference, model evaluation.}

\project{\href{https://github.com/WaseemRaza844/agentic-ai-projects-by-was}{Agentic AI and RAG Applications}}{Portfolio Project}{2026}
Built Python workflows that combine retrieval, vector search, tool-using agents, state/memory, evaluation, and multi-step orchestration. Used iterative agent-assisted development and validation to structure implementation, debug workflows, and verify expected behavior across multi-stage tasks.\\[-1pt]
{\small\textit{Tools \& methods:} Python, LangGraph, LangChain, RAG, vector databases, embeddings, tool calling, agent state/memory, workflow evaluation, coding-agent-assisted iteration.}

\section{Selected Technical Training}
\cert{\href{https://www.coursera.org/account/accomplishments/verify/VU3Y5TBLVLZT}{Machine Learning}}{Stanford Online / Coursera | Completed 03/2020}{Regression, classification, regularization, clustering, model evaluation, bias/variance, and practical model-selection workflows for supervised and unsupervised learning.}
\cert{\href{https://www.coursera.org/account/accomplishments/verify/SG4GDJ4CLPVL}{Generative AI Engineering and Fine-Tuning Transformers}}{IBM / Coursera | Completed 08/2026}{Transformer architectures, pretrained models, task adaptation, fine-tuning strategies, tokenization concepts, training/inference tradeoffs, and LLM evaluation.}
\cert{\href{https://www.coursera.org/account/accomplishments/verify/Y6NP4GWLS4DZ}{Introduction to Deep Learning \& Neural Networks with Keras}}{IBM / Coursera | Completed 08/2026}{Neural-network architecture, activation functions, optimization, model training, validation, and iterative model improvement using Keras-based workflows.}
\cert{\href{https://www.coursera.org/account/accomplishments/verify/WRA89JEJU8EE}{Python for Data Science, AI \& Development}}{IBM / Coursera | Completed 07/2026}{Python programming, data structures, data manipulation, NumPy/Pandas-oriented analytical workflows, and reusable application development for data science and AI.}
\cert{\href{https://www.coursera.org/account/accomplishments/verify/R6248BNQAC0D}{Introduction to Computer Vision with TensorFlow}}{Google Cloud / Coursera | Completed 08/2024}{Image preprocessing, convolutional-network concepts, image classification, model construction, training, validation, and practical evaluation with TensorFlow workflows.}

\section{Selected Research \& Validation Evidence}
\pub{\href{https://doi.org/10.1109/PIMRC59610.2024.10817403}{An AI-Driven Framework for Enhancing Resilience in Propagation Models to Enable Digital Twin}}{IEEE PIMRC, 2024}{Robust model development and evaluation under scarce measured data and distribution shift using neural pretraining and generative augmentation.}
\pub{Real-World Medication Adherence Monitoring Using Hybrid Neural Networks and Smartwatches}{HealthINF, Accepted 2026}{Deep-learning sequence classification, multimodal wearable signals, cross-user evaluation, and real-world reliability analysis.}
\pub{Data-Driven Intelligent Outage Management for High Shadowing Environments in 5G\&B Networks}{IEEE SmartNets, 2024}{Machine-learning-based anomaly/outage detection, data-driven decision support, and validation under challenging operational conditions.}
\pub{\href{https://doi.org/10.3390/s25247651}{On Multi-Parameter Optimization and Proactive Reliability in 5G and Beyond Cellular Networks}}{Sensors, 2025}{Predictive modeling, optimization, reliability analysis, and systematic evaluation for data-driven network operation.}

\section{Coding, Review \& Reproducibility Practices}
\begin{itemize}
\item Build Python workflows as testable, repeatable stages for data preparation, model execution, evaluation, error analysis, and result validation rather than ad-hoc notebook-only experiments.
\item Use Git, command-line tooling, automated tests, GitHub Actions, and coding-agent-assisted iteration to review implementations, debug failures, refine code, and verify objective correctness.
\item Comfortable reviewing ML code and experimental outputs for data leakage, shape/type errors, metric misuse, reproducibility gaps, cross-domain failure modes, and inconsistencies between intended and actual model behavior.
\item Evaluate candidate implementations against explicit correctness and performance criteria, documenting assumptions, metrics, failure cases, and reproducibility before accepting results.
\end{itemize}

\end{document}
